% (c) Nikita Lisitsa, lisyarus@gmail.com, 2026

\documentclass[10pt]{beamer}

\usepackage[T2A]{fontenc}
\usepackage[russian]{babel}
\usepackage{minted}

\usepackage[most]{tcolorbox}
\tcbuselibrary{minted}

\usepackage{graphicx}
\graphicspath{ {./images/} }

\usepackage{adjustbox}

\usepackage{color}
\usepackage{soul}

\usepackage{hyperref}

\usetheme{metropolis}

\definecolor{red}{rgb}{1,0,0}
\definecolor{codebg}{RGB}{29,35,49}
\setminted{fontsize=\footnotesize}

\makeatletter
\newcommand{\slideimage}[1]{
  \begin{figure}
    \begin{adjustbox}{width=\textwidth, totalheight=\textheight-2\baselineskip-2\baselineskip,keepaspectratio}
      \includegraphics{#1}
    \end{adjustbox}
  \end{figure}
}
\makeatother

\title{Язык программирования C++}
\subtitle{Лекция 9: ссылки, перегрузка функций, new/delete, ООП, классы, инкапсуляция}
\date{2026}
\author{Никита Лисица}

\setbeamertemplate{footline}[frame number]

\begin{document}

\frame{\titlepage}

\begin{frame}
\frametitle{Немного истории}
\begin{itemize}
\item \textbf{1979--1985}: Бьёрн Страуструп, работая в Bell Labs, проектирует и реализует новый язык на основе C
\pause
\item Конец \textbf{80-х -- 90-е}: Язык набирает популярность
\pause
\item \textbf{1998}: Первый официальный стандарт -- C++98
\pause
\item \textbf{2003}: Дополнение к стандарту -- C++03
\pause
\item \textbf{2000-е}: Язык развивается медленно в сравнении с другими языками того же времени (Java, C\#, Python, Scala, F\#, ...)
\pause
\item \textbf{2011}: Стандарт C++11, привнёсший много нового и изменивший взгляд на C++ (т.н. \textit{Modern C++})
\pause
\item \textbf{Дальше}: Новый стандарт каждые 3 года (C++14, C++17, C++21, C++23, C++26)
\end{itemize}
\end{frame}

\begin{frame}
\frametitle{Немного контекста}
\begin{itemize}
\item Начиная с 90х годов, C++ -- \textbf{основной язык} для низкоуровневых, высокопроизводительных систем
\pause
\item Все браузеры, программы для обработки графики и видео, крупные игровые движки и т.д. написаны на C++
\pause
\item Сегодня C++ \textit{частично} замещается другими языками (Rust, Go, D)
\pause
\item Тем не менее, согласно \href{https://www.tiobe.com/tiobe-index/}{\alert{\underline{TIOBE Index}}}, C++ всё ещё на \textbf{третьем месте} по популярности (после Python и C)
\end{itemize}
\end{frame}

\begin{frame}
\frametitle{Немного контекста}
\begin{itemize}
\item С++ -- безумно сложный язык (в стандарте сейчас около 2400 страниц)
\pause
\item Пытается совместить низкоуровневость C и высокоуровневые концепции
\pause
\item Взаимодействие всех аспектов языка приводит к тысячам нюансов и неочевидных спецэффектов
\pause
\item В балансе между простотой использования и мощностью языковых конструкций C++ выбирает второе :)
\pause
\item Полезно изучить, даже если не планируете его использовать
\end{itemize}
\end{frame}

\begin{frame}[fragile]
\frametitle{Процесс компиляции}
\begin{itemize}
\usemintedstyle{tango}
\item Компиляция C++ происходит так же, как компиляция C
\pause
\item Исходный код делится на файлы, которые компилируются по отдельности, и собираются вместе линковщиком в результирующую программу/библиотеку
\pause
\item Файлы с кодом на C++ обычно имеют расширение \nolinkurl{.cpp}, реже \nolinkurl{.cxx} или \nolinkurl{.cc}
\pause
\item Заголовочные файлы на C++ обычно имеют расширение \nolinkurl{.h} или \nolinkurl{.hpp} (чтобы отличать от C-файлов), реже \nolinkurl{.hh} или \nolinkurl{.hxx}
\pause
\item \textbf{NB:} Заголовочные файлы стандартной библиотеки C++ не имеют расширения, напр. \mintinline{cpp}|#include <thread>|
\end{itemize}
\end{frame}

\begin{frame}[fragile]
\frametitle{C++ vs C}
\begin{itemize}
\item C++ наследует почти всё, что есть в C, но есть отличия
\pause
\item Самое важное -- нет неявных приведений типов указателей:
\usemintedstyle{lightbulb}
\begin{tcblisting}{listing engine=minted, minted language=cpp, minted options={fontsize=\scriptsize}, listing only, colback=codebg, colframe=codebg, rounded corners}
int *pi = ...;
char *pc = ...;

// Можно в C, ошибка компиляции в C++
pi = pc;

// В C++ нужно явное приведение типа
pi = (int*)pc;

// Приводить к void* в C++ всё ещё можно
void *pv = pi;

// Но не обратно (ошибка компиляции в C++)
pc = pv;
\end{tcblisting}
\end{itemize}
\end{frame}

\begin{frame}[fragile]
\frametitle{Выделение/освобождение памяти: new/delete}
\begin{itemize}
\usemintedstyle{tango}
\item В C++ обычно память выделяют/освобождают не \textit{функциями} \mintinline{cpp}|malloc/free|, а \textit{операторами} \mintinline{cpp}|new/delete|:
\usemintedstyle{lightbulb}
\begin{tcblisting}{listing engine=minted, minted language=cpp, minted options={fontsize=\scriptsize}, listing only, colback=codebg, colframe=codebg, rounded corners}
// Аналог malloc(sizeof(int))
int *pi = new int;

// Аналог free(pi)
delete pi;

// Аналог malloc(N * sizeof(int))
int *array = new int[N];

// Аналог free(array)
delete [] array;
\end{tcblisting}
\pause
\usemintedstyle{tango}
\item \mintinline{cpp}|new| и \mintinline{cpp}|delete| -- \textit{операторы} языка, они обрабатываются напрямую компилятором, и знают про тип, который выделяется
\pause
\item Это важно для \mintinline{cpp}|sizeof| и для конструкторов/декструкторов (о них позже)
\end{itemize}
\end{frame}

\begin{frame}[fragile]
\frametitle{Выделение/освобождение памяти: new/delete}
\begin{itemize}
\usemintedstyle{tango}
\item Указатели, выделенные с помощью \mintinline{cpp}|malloc|, \textbf{нужно} освобождать с помощью \mintinline{cpp}|free|
\pause
\item Указатели, выделенные с помощью \mintinline{cpp}|new|, \textbf{нужно} освобождать с помощью \mintinline{cpp}|delete|
\pause
\item Указатели, выделенные с помощью \mintinline{cpp}|new[]|, \textbf{нужно} освобождать с помощью \mintinline{cpp}|delete[]|
\pause
\item Все три способа могут (и будут) внутри работать по-разному (например, запоминать количество объектов в \mintinline{cpp}|new[]|), и быть несовместимыми
\end{itemize}
\end{frame}

\begin{frame}[fragile]
\frametitle{Ссылки (references)}
\begin{itemize}
\usemintedstyle{tango}
\item Указатели -- один из основных инструментов в C, но бывают неудобны:
\pause
\begin{itemize}
\item Нужно помнить про взятие адреса \mintinline{cpp}|&x| и разыменовывание \mintinline{cpp}|*p|
\pause
\item Могут быть нулевыми \mintinline{cpp}|int * p = NULL;|
\pause
\item Адрес, на который они указывают, может меняться \mintinline{cpp}|p = &x|
\end{itemize}
\pause
\item В C++ есть альтернатива: \textbf{ссылки}
\end{itemize}
\end{frame}

\begin{frame}[fragile]
\frametitle{Ссылки (references)}
\begin{itemize}
\usemintedstyle{tango}
\item Тип ссылки на \mintinline{cpp}|T| обозначается как \mintinline{cpp}|T&| (например, \mintinline{cpp}|int&| -- ссылка на \mintinline{cpp}|int|)
\pause
\usemintedstyle{lightbulb}
\begin{tcblisting}{listing engine=minted, minted language=cpp, minted options={fontsize=\scriptsize}, listing only, colback=codebg, colframe=codebg, rounded corners}
int x = ...;
// Здесь & - не взятие адреса, а часть типа
// int& означает ссылку на int
int &ref = x;

// Ссылку обязательно нужно проинициализировать!
int &ref2; // Ошибка компиляции

// Ссылку нельзя перенаправить на другой объект
int y = 42;
ref = y; // присваивает переменной x значение 42
// ref всё ещё указывает на x
\end{tcblisting}
\end{itemize}
\end{frame}

\begin{frame}[fragile]
\frametitle{Ссылки (references)}
\usemintedstyle{lightbulb}
\begin{tcblisting}{listing engine=minted, minted language=cpp, minted options={fontsize=\scriptsize}, listing only, colback=codebg, colframe=codebg, rounded corners}
// Разыменование указателя возвращает ссылку
int *p = ...;
int &ref = *p;

// У ссылки можно взять адрес - получится
// адрес объекта, на который она ссылается
int *q = &ref;
assert(p == q);
\end{tcblisting}
\end{frame}

\begin{frame}[fragile]
\frametitle{Ссылки (references)}
\begin{itemize}
\usemintedstyle{tango}
\item В целом, ссылка на объект ведёт себя как \textit{синоним} имени этого объекта
\pause
\item Самый близкий аналог в C -- \mintinline{cpp}|T *const|, но без операторов взятия адреса и разыменования
\pause
\item В бинарном (скомипилированном) коде ссылка выглядит так же, как указатель -- просто адрес в памяти
\end{itemize}
\end{frame}

\begin{frame}[fragile]
\frametitle{Ссылки (references)}
\begin{itemize}
\usemintedstyle{tango}
\item Ссылка может указывать на константный или неконстантный объект:
\usemintedstyle{lightbulb}
\begin{tcblisting}{listing engine=minted, minted language=cpp, minted options={fontsize=\scriptsize}, listing only, colback=codebg, colframe=codebg, rounded corners}
int x = ...;
int &ref = x;
const int &cref = x;

// Через неконстантную ссылку можно
// менять значение переменной x
ref = 42;

// Через константную ссылку нельзя - 
// ошибка компиляции
cref = 42;

const int y = ...;
// Ошибка компиляции: нельзя взять неконстантную
// ссылку на константную переменную
int &yref = y;
\end{tcblisting}
\end{itemize}
\end{frame}

\begin{frame}[fragile]
\frametitle{Ссылки (references)}
\begin{itemize}
\usemintedstyle{tango}
\item Аргументы функции могут быть ссылками:
\usemintedstyle{lightbulb}
\begin{tcblisting}{listing engine=minted, minted language=cpp, minted options={fontsize=\scriptsize}, listing only, colback=codebg, colframe=codebg, rounded corners}
void swap(int &x, int &y)
{
  int temp = x;
  x = y;
  y = temp;
}

// Для вызова функции swap не нужны
// операторы взятия адреса:
int a = 3, b = 5;
swap(a, b);
\end{tcblisting}
\pause
\item У этого есть минус: по вызову функции не всегда понятно, что копируется, а что берётся по ссылке (нужно смотреть сигнатуру вызываемой функции)
\end{itemize}
\end{frame}

\begin{frame}[fragile]
\frametitle{Перегрузка функций (function overloading)}
\begin{itemize}
\usemintedstyle{tango}
\item В языке C нельзя иметь несколько функций с одним именем (т.е. \textit{перегрузить} имя функции)
\pause
\item В C++ -- \textbf{можно}! Но они должны отличаться набором (количеством и/или типами) аргументов:
\usemintedstyle{lightbulb}
\begin{tcblisting}{listing engine=minted, minted language=cpp, minted options={fontsize=\scriptsize}, listing only, colback=codebg, colframe=codebg, rounded corners}
int max(int x, int y) { ... }

int max(int x, int y, int z) { ... }

// Ошибка компиляции: отличается только тип
// возвращаемого значения
double max(int x, int y) { ... }

// OK - отличается тип первого аргумента
double max(double x, int y) { ... }
\end{tcblisting}
\pause
\item Очень полезно, когда нужно реализовать похожую логику для разных типов данных или с разными настройками
\end{itemize}
\end{frame}

\begin{frame}[fragile]
\frametitle{Перегрузка функций: name mangling}
\begin{itemize}
\usemintedstyle{tango}
\item После компиляции линковщик подставляет адреса нужных функций в места их вызова
\pause
\item Если функция перегружена, есть несколько функций с одинаковым именем -- какую он должен выбрать?
\pause
\item Можно было бы научить линковщик правилам перегрузки C++, но не хотелось бы
\pause
\item Решение: \textbf{name mangling}
\end{itemize}
\end{frame}

\begin{frame}[fragile]
\frametitle{Перегрузка функций: name mangling}
\begin{itemize}
\usemintedstyle{tango}
\item При компиляции C++ имена функций преображаются компилятором так, чтобы функции с разным количеством/типами аргументов превращались в разные \textit{mangled} имена
\pause
\item Конкретный алгоритм преобразования зависит от системы и компилятора (т.е. является частью ABI)
\pause
\item Name mangling важен для классов и неймспейсов, о них позже
\end{itemize}
\end{frame}

\begin{frame}[fragile]
\frametitle{Перегрузка функций: name mangling}
\begin{itemize}
\usemintedstyle{tango}
\item Под Linux можно узнать mangled имена командой \mintinline{bash}|readelf -s file.o|:
\usemintedstyle{lightbulb}
\begin{tcblisting}{listing engine=minted, minted language=cpp, minted options={fontsize=\scriptsize}, listing only, colback=codebg, colframe=codebg, rounded corners}
// -> _Z3maxii
int max(int x, int y) { ... }

// -> _Z3maxiii
int max(int x, int y, int z) { ... }

// -> _Z3maxff
double max(double x, int y) { ... }

// -> _Z3maxPcS_
char * max(char *p, char * q) { ... }
\end{tcblisting}
\end{itemize}
\end{frame}

\begin{frame}[fragile]
\frametitle{ООП (объектно-ориентированное программирование)}
\begin{itemize}
\usemintedstyle{tango}  
\item Обычно говорят, что C++ -- объектно-ориентированный язык
\pause
\item Правильнее сказать, что C++ поддерживает объектно-ориентированную парадигму (методологию, стиль) программирования
\pause
\item Можно писать на C++, вообще не используя ООП-механизмы, но обычно так не делают :)
\end{itemize}
\end{frame}

\begin{frame}[fragile]
\frametitle{ООП (объектно-ориентированное программирование)}
\begin{itemize}
\usemintedstyle{tango}
\item Что вообще такое ООП? \pause Три принципа:
\pause
\item Инкапсуляция
\begin{itemize}
\item Механизм языка, ограничивающий доступ одних компонентов программы к другим
\item Языковая конструкция, связывающая данные с методами для их обработки
\end{itemize}
\pause
\item Наследование
\begin{itemize}
\item Механизм, позволяющий `отнаследовать' данные и функциональность одного объекта другим, способствуя переиспользованию кода
\end{itemize}
\pause
\item Полиморфизм
\begin{itemize}
\item Способность одного и того же кода обрабатывать разные типы данных
\end{itemize}
\end{itemize}
\end{frame}

\begin{frame}[fragile]
\frametitle{ООП в C++: Классы}
\begin{itemize}
\usemintedstyle{tango}
\item В C++ объектно-ориентированное программирование выражается посредством \textbf{классов}
\pause
\item Проще всего смотреть на классы как на навороченные структуры (точное отличие классов и структур обсудим позднее)
\pause
\item Объявление класса:
\usemintedstyle{lightbulb}
\begin{tcblisting}{listing engine=minted, minted language=cpp, minted options={fontsize=\scriptsize}, listing only, colback=codebg, colframe=codebg, rounded corners}
// Array - имя типа, аналогично структурам в С
class Array {
  // Поля (данные) класса, аналогично структурам в C
  int *data;
  size_t size;

  // Методы: функции, которые можно вызывать у объектов класса
  void set(size_t index, int value);
  void fill(int value);
};
\end{tcblisting}
\end{itemize}
\end{frame}

\begin{frame}[fragile]
\frametitle{Динамический массив без классов}
\begin{itemize}
\usemintedstyle{lightbulb}
\begin{tcblisting}{listing engine=minted, minted language=cpp, minted options={fontsize=\scriptsize}, listing only, colback=codebg, colframe=codebg, rounded corners}
int *array = new int[size];

// Как-то работаем с массивом
for (...)
  array[i] = ...;

delete[] array;
\end{tcblisting}
\pause
\item Чем плох такой динамический массив?
\pause
\usemintedstyle{tango}
\begin{itemize}
\item Нужно не забыть выделить память
\item Нужно не выйти за его границы
\item Нужно не перепутать размер (напр. в цикле \mintinline{cpp}|for|)
\item Нужно не забыть освободить память
\end{itemize}
\end{itemize}
\end{frame}

\begin{frame}[fragile]
\frametitle{Класс динамического массива}
\usemintedstyle{lightbulb}
\begin{tcblisting}{listing engine=minted, minted language=cpp, minted options={fontsize=\scriptsize}, listing only, colback=codebg, colframe=codebg, rounded corners}
// array.h

class Array {
// Инкапсуляция: приватные поля недоступны снаружи объекта
private:
  size_t size;
  int *data;

// Публичный интерфейс класса
public:
  
  // Конструктор - создаёт объект
  Array(size_t size);

  // Деструктор - удаляет объект
  ~Array();

  // Доступ к элементам
  int get(size_t index);
  void set(size_t index, value);
};
\end{tcblisting}
\end{frame}

\begin{frame}[fragile]
\frametitle{Класс динамического массива}
\usemintedstyle{lightbulb}
\begin{tcblisting}{listing engine=minted, minted language=cpp, minted options={fontsize=\scriptsize}, listing only, colback=codebg, colframe=codebg, rounded corners}
// array.сpp

// Реализация конструктора
Array::Array(size_t size) {
{
  // this - указатель на объект, у которого вызвали метод
  this->size = size;
  data = new int[size];
}

// Реализация деструктора
Array::~Array() {
  delete[] data;
}

// Реализация методов
int Array::get(size_t index) {
  return data[index];
}

void Array::set(size_t index, value) {
  data[index] = value;
}
\end{tcblisting}
\end{frame}

\begin{frame}[fragile]
\frametitle{Класс динамического массива}
\usemintedstyle{lightbulb}
\begin{tcblisting}{listing engine=minted, minted language=cpp, minted options={fontsize=\scriptsize}, listing only, colback=codebg, colframe=codebg, rounded corners}
#include <array.h>

void test_array() {
  // Вызов конструктора, size = 100
  Array a(100);

  // Ошибка компиляции: нет параметра конструктора
  Array b;

  // Ошибка компиляции: data - приватное поле
  a.data[10] = 42;

  // OK, вызов метода класса
  a.set(10, 42);

  // Здесь автоматически(!) вызовется деструктор a
  // Как будто мы написали a.~Array() (не надо так писать самим!)
}
\end{tcblisting}
\end{frame}

\begin{frame}[fragile]
\frametitle{Немного о терминологии}
\usemintedstyle{tango}
\begin{itemize}
\item \mintinline{cpp}|int x;|
\pause
\begin{itemize}
\item \mintinline{cpp}|int| -- \textbf{тип} переменной
\item \mintinline{cpp}|x| -- \textbf{имя} переменной
\end{itemize}
\pause
\item \mintinline{cpp}|Array a;|
\pause
\begin{itemize}
\item \mintinline{cpp}|Array| -- \textbf{класс (class)} (он в том числе является типом!)
\item \mintinline{cpp}|a| -- \textbf{объект/экземпляр (object/instance)} класса (во многом ведёт себя как переменная)
\item \mintinline{cpp}|data| -- \textbf{поле (field/member)} класса
\item \mintinline{cpp}|set()| -- \textbf{метод (method)} класса
\end{itemize}
\end{itemize}
\end{frame}

\begin{frame}[fragile]
\frametitle{Класс динамического массива}
\usemintedstyle{lightbulb}
\begin{itemize}
\item Добавим \textit{метод}, позволяющий получить размер массива:
\pause
\begin{tcblisting}{listing engine=minted, minted language=cpp, minted options={fontsize=\scriptsize}, listing only, colback=codebg, colframe=codebg, rounded corners}
class Array {
...
public:
  ...

  // Методы можно реализовывать прямо в объявлении класса!
  // Тогда они автоматически считаются inline
  size_t get_size() {
    return size;
  }
};
\end{tcblisting}
\pause
\usemintedstyle{tango}
\item Пришлось назвать его \mintinline{cpp}|get_size()| а не \mintinline{cpp}|size()|, потому что такое имя конфликтовало бы с \textit{полем} \mintinline{cpp}|size|
\pause
\item Из-за этого приватные поля часто называют с префиксами/суффиксами, например \mintinline{cpp}|mSize| (m от member) или \mintinline{cpp}|size_|
\end{itemize}
\end{frame}

\begin{frame}[fragile]
\frametitle{Конструкторы}
\usemintedstyle{lightbulb}
\begin{itemize}
\item Конструктор класса -- просто особенная функция (но без возвращаемого значения!), и его можно перегружать:
\begin{tcblisting}{listing engine=minted, minted language=cpp, minted options={fontsize=\scriptsize}, listing only, colback=codebg, colframe=codebg, rounded corners}
// array.h

class Array {
...
public:
  // Выделить память, оставить непроинициализированной
  Array(size_t size);

  // Выделить память, заполнить значением value
  Array(size_t size, int value);

  // Создать массив нулевого размера
  Array();
};
\end{tcblisting}
\end{itemize}
\end{frame}

\begin{frame}[fragile]
\frametitle{Конструкторы}
\usemintedstyle{lightbulb}
\begin{tcblisting}{listing engine=minted, minted language=cpp, minted options={fontsize=\scriptsize}, listing only, colback=codebg, colframe=codebg, rounded corners}
// array.cpp

Array::Array(size_t size) {
  this->size = size;
  data = new int[size];
}

Array::Array(size_t size, int value) {
  this->size = size;
  data = new int[size];
  for (size_t i = 0; i < size; ++i)
    data[i] = value;
}

Array::Array() {
  size = 0;
  data = nullptr; // более безопасный аналог NULL
}
\end{tcblisting}
\end{frame}

\begin{frame}[fragile]
\frametitle{Конструктор по умолчанию}
\usemintedstyle{lightbulb}
\begin{itemize}
\usemintedstyle{tango}
\item Конструктор без аргументов \mintinline{cpp}|Array()| называется конструктором по умолчанию (default constructor)
\pause
\usemintedstyle{lightbulb}
\item Его нужно вызывать без скобок:
\begin{tcblisting}{listing engine=minted, minted language=cpp, minted options={fontsize=\scriptsize}, listing only, colback=codebg, colframe=codebg, rounded corners}
// OK, вызов конструктора по умолчанию
Array a;
\end{tcblisting}
\pause
\item А со скобками не получится:
\begin{tcblisting}{listing engine=minted, minted language=cpp, minted options={fontsize=\scriptsize}, listing only, colback=codebg, colframe=codebg, rounded corners}
// Не OK, это объявление функции a() без аргументов,
// возвращающей объект класса Array
Array a();
\end{tcblisting}
\end{itemize}
\end{frame}

\begin{frame}[fragile]
\frametitle{Неявный конструктор по умолчанию}
\begin{itemize}
\usemintedstyle{lightbulb}
\item Если не написать ни одного конструктора, компилятор сам добавит конструктор по умолчанию, который ничего не делает:
\begin{tcblisting}{listing engine=minted, minted language=cpp, minted options={fontsize=\scriptsize}, listing only, colback=codebg, colframe=codebg, rounded corners}
Array::Array() {}
\end{tcblisting}
\pause
\item Аналогично с деструктором:
\begin{tcblisting}{listing engine=minted, minted language=cpp, minted options={fontsize=\scriptsize}, listing only, colback=codebg, colframe=codebg, rounded corners}
Array::~Array() {}
\end{tcblisting}
\pause
\item \textbf{NB:} конструкторы/деструкторы класса всегда автоматически вызывают конструкторы/деструкторы всех полей, об этом позже
\end{itemize}
\end{frame}

\begin{frame}[fragile]
\frametitle{C-style инкапсуляция}
\usemintedstyle{lightbulb}
\begin{tcblisting}{listing engine=minted, minted language=cpp, minted options={fontsize=\scriptsize}, listing only, colback=codebg, colframe=codebg, rounded corners}
typedef struct {
  size_t size;
  int *data;
} array_t;

void init(array_t *array, size_t size) {
  array->size = size;
  array->data = malloc(size * sizeof(int));
}

void destroy(array_t *array) {
  free(array->data);
}

int main() {
  array_t a;
  init(&a, 100);
  destroy(&a);
}
\end{tcblisting}
\end{frame}

\begin{frame}[fragile]
\frametitle{Релизация инкапсуляции в C++}
\usemintedstyle{lightbulb}
\begin{itemize}
\usemintedstyle{tango}
\item Компилятор превращает код с классами во что-то похожее на C-style инкапсуляцию
\pause
\item Все проверки на public/private происходят во время компиляции, в бинарном коде они уже не нужны
\pause
\item Методы класса превращаются в функции, принимающие дополнительный параметр \mintinline{cpp}|this| (часто у таких функций другой calling convention, см. thiscall)
\pause
\item Конструкторы и деструкторы, аналогично методам, превращаются в функции, и вызываются компилятором автоматически в нужных местах
\end{itemize}
\end{frame}

\end{document}